
\subsubsection{2.1 Spurious Correlations and Group Robustness}

Sagawa et al. (2020) introduced Group DRO, which minimizes worst-case loss over predefined groups with strong regularization, achieving substantial improvements on Waterbirds and CelebA benchmarks. However, Group DRO requires group annotations during training.

Just Train Twice \citep{liu2021just} trains an initial ERM model for few epochs, identifies misclassified examples (which correlate with minority groups), and upweights them in a second training run. Learning from Failure \citep{nam2020learning} explicitly trains a biased network and uses its failures to guide debiasing. Both methods exploit the observation that spurious features are learned early but require two-stage training.

Deep Feature Reweighting \citep{kirichenko2022last} demonstrates that pretrained ERM features are sufficient for strong worst-group accuracy when the last layer is retrained on group-balanced data. This finding directly relates to the present negative result: if pretrained features already separate spurious and core concepts, there may be no emergence dynamics left to observe.

\subsubsection{2.2 Simplicity Bias and Learning Dynamics}

Shah et al. (2020) established that SGD exhibits simplicity bias: networks learn the simplest predictive features first and may never learn more complex features even when they have higher predictive power. Gradient Starvation \citep{pezeshki2021gradient} provides a dynamical systems perspective: cross-entropy minimization on features with different frequencies causes some features to receive diminishing gradient signal.

The present work extends this literature by testing whether emergence uniformity (variance across sample subsets) differs between spurious and core features. The negative result suggests that while emergence timing may differ during training, this signal is not recoverable from pretrained representations.

\subsubsection{2.3 Linear Probing for Representation Analysis}

Linear probes are standard tools for analyzing pretrained representations \citep{alain2017understanding}. The CLIP evaluation protocol \citep{radford2021learning} uses logistic regression with regularization-strength sweeps to assess representation quality.

The present experiment adapted this protocol by treating regularization strength (C) as an epoch proxy and computing CV across sample subsets. The failure reveals a category error: C-sweep produces representation quality scores, not learning trajectory data. Convex optimization converges in a single pass regardless of C; there is no emergence to measure.

